% GENERATED FILE. Edit canonical lessons, then run npm run generate:books.
% canonical-source-hash: fnv1a64:028a41c981f345ba
% canonical-lessons: UR-C468-mahak, UR-C468-itr, UR-C468-lams, UR-C468-khujana, UR-C468-kharish

\chapter{Smells and Itches}
\label{ch:ur-smells-and-itches}
\hlchaptermodality{468}

\begin{chapteropening}
\textbf{By the end of this chapter:} \emph{I can talk about smells and itches.}

Complete the last lesson of chapter 468: I can talk about smells and itches.
\end{chapteropening}

\section[\texorpdfstring{\ur{مہک} (mahak)}{مہک (mahak)}]{\ur{مہک} (mahak) — a scent}

\label{lesson:UR-C468-mahak}



\begin{quote}
Before the new one: say the Urdu for to suck, then the Urdu for to spit.
\end{quote}

\subsection*{You'll want to know: \ur{مہک}}
\textbf{\ur{مہک}} — \emph{mahak} — \textquotedblleft{}a scent\textquotedblright{}.

\begin{quote}
\ur{کمرے} \ur{میں} \ur{پھولوں} \ur{کی} \ur{مہک} \ur{ہے۔}
\end{quote}

\begin{quote}
\emph{kamre meṅ phūloṅ kī mahak hai.}
\end{quote}

\begin{quote}
There is a scent of flowers in the room.
\end{quote}

Feminine.

\begin{cousinweb}
\textbf{\ur{مہک}} (\emph{mahak}) is inherited: CDIAL lists the Hindi cousin \emph{mahak}, \textquotedblleft{}fragrance\textquotedblright{}, under Sanskrit \emph{magha} (CDIAL 9703).
\end{cousinweb}

\begin{grammarlens}[title={What you've built}]
One more word for describing what you lived and felt.
\end{grammarlens}

\subsection*{Your turn}
\begin{itemize}
\raggedright
  \item \emph{Say it:} \emph{mahak}
  \item \emph{Say it:} \emph{mahak}, once more
  \item \emph{Say it:} \emph{kamre meṅ phūloṅ kī mahak hai.} — There is a scent of flowers in the room.
  \item \emph{Recall:} say the Urdu for to suck, then the Urdu for to spit, then say \emph{mahak} again
\end{itemize}

\begin{tcolorbox}[breakable,colback=teal!4,colframe=teal!35!black,title={Before you move on}]
What does \ur{مہک} mean? (\textquotedblleft{}A scent\textquotedblright{}.) Say it once more.
\end{tcolorbox}
\section[\texorpdfstring{\ur{عطر} (itr)}{عطر (itr)}]{\ur{عطر} (itr) — perfume; attar}

\label{lesson:UR-C468-itr}



\begin{quote}
Before the new one: say the Urdu for to spit, then the Urdu for a scent.
\end{quote}

\subsection*{You'll want to know: \ur{عطر}}
\textbf{\ur{عطر}} — \emph{itr} — \textquotedblleft{}perfume; attar\textquotedblright{}.

\begin{quote}
\ur{اس} \ur{عطر} \ur{کی} \ur{خوشبو} \ur{بڑی} \ur{اچھی} \ur{ہے۔}
\end{quote}

\begin{quote}
\emph{is itr kī khushbū baṛī achchhī hai.}
\end{quote}

\begin{quote}
This perfume smells lovely.
\end{quote}

Masculine.

\begin{cousinweb}
\textbf{\ur{عطر}} (\emph{itr}) is usually traced, through Persian, to Arabic \emph{ʿiṭr}, \textquotedblleft{}perfume\textquotedblright{}.
\end{cousinweb}

\begin{grammarlens}[title={What you've built}]
One more word for describing what you lived and felt.
\end{grammarlens}

\subsection*{Your turn}
\begin{itemize}
\raggedright
  \item \emph{Say it:} \emph{itr}
  \item \emph{Say it:} \emph{itr}, once more
  \item \emph{Say it:} \emph{is itr kī khushbū baṛī achchhī hai.} — This perfume smells lovely.
  \item \emph{Recall:} say the Urdu for to spit, then the Urdu for a scent, then say \emph{itr} again
\end{itemize}

\begin{tcolorbox}[breakable,colback=teal!4,colframe=teal!35!black,title={Before you move on}]
What does \ur{عطر} mean? (\textquotedblleft{}Perfume; attar\textquotedblright{}.) Say it once more.
\end{tcolorbox}
\section[\texorpdfstring{\ur{لمس} (lams)}{لمس (lams)}]{\ur{لمس} (lams) — the sense of touch}

\label{lesson:UR-C468-lams}



\begin{quote}
Before the new one: say the Urdu for a scent, then the Urdu for perfume.
\end{quote}

\subsection*{You'll want to know: \ur{لمس}}
\textbf{\ur{لمس}} — \emph{lams} — \textquotedblleft{}the sense of touch\textquotedblright{}.

\begin{quote}
\ur{ماں} \ur{کے} \ur{لمس} \ur{سے} \ur{بچہ} \ur{سو} \ur{گیا۔}
\end{quote}

\begin{quote}
\emph{māṅ ke lams se bachchā so gayā.}
\end{quote}

\begin{quote}
At his mother's touch the baby fell asleep.
\end{quote}

Masculine.

\begin{cousinweb}
\textbf{\ur{لمس}} (\emph{lams}) is usually traced, through Persian, to Arabic \emph{lams}, \textquotedblleft{}touching\textquotedblright{}.
\end{cousinweb}

\begin{grammarlens}[title={What you've built}]
One more word for describing what you lived and felt.
\end{grammarlens}

\subsection*{Your turn}
\begin{itemize}
\raggedright
  \item \emph{Say it:} \emph{lams}
  \item \emph{Say it:} \emph{lams}, once more
  \item \emph{Say it:} \emph{māṅ ke lams se bachchā so gayā.} — At his mother's touch the baby fell asleep.
  \item \emph{Recall:} say the Urdu for a scent, then the Urdu for perfume, then say \emph{lams} again
\end{itemize}

\begin{tcolorbox}[breakable,colback=teal!4,colframe=teal!35!black,title={Before you move on}]
What does \ur{لمس} mean? (\textquotedblleft{}The sense of touch\textquotedblright{}.) Say it once more.
\end{tcolorbox}
\section[\texorpdfstring{\ur{کھجانا} (khujānā)}{کھجانا (khujānā)}]{\ur{کھجانا} (khujānā) — to scratch an itch}

\label{lesson:UR-C468-khujana}



\begin{quote}
Before the new one: say the Urdu for perfume, then the Urdu for the sense of touch.
\end{quote}

\subsection*{You'll want to know: \ur{کھجانا}}
\textbf{\ur{کھجانا}} — \emph{khujānā} — \textquotedblleft{}to scratch an itch\textquotedblright{}.

\begin{quote}
\ur{مچھر} \ur{کے} \ur{کاٹنے} \ur{پر} \ur{وہ} \ur{ہاتھ} \ur{کھجا} \ur{رہا} \ur{تھا۔}
\end{quote}

\begin{quote}
\emph{machchhar ke kāṭne par voh hāth khujā rahā thā.}
\end{quote}

\begin{quote}
He was scratching his hand where the mosquito had bitten.
\end{quote}

Past, with \ur{نے} (he or she …): \ur{اس} \ur{نے} \ur{سر} \ur{کھجایا}, \emph{us ne sar khujāyā}.

\begin{cousinweb}
\textbf{\ur{کھجانا}} (\emph{khujānā}) is inherited: CDIAL lists the Hindi cousin \emph{khujānā}, \textquotedblleft{}to itch\textquotedblright{}, under Sanskrit \emph{kharju}, \textquotedblleft{}itching, scratching\textquotedblright{} (CDIAL 3827); \textbf{\ur{کھجلی}}, \textquotedblleft{}an itch\textquotedblright{}, is from the same word.
\end{cousinweb}

\begin{grammarlens}[title={What you've built}]
One more word for describing what you lived and felt.
\end{grammarlens}

\subsection*{Your turn}
\begin{itemize}
\raggedright
  \item \emph{Say it:} \emph{khujānā}
  \item \emph{Say it:} \emph{khujānā}, once more
  \item \emph{Say it:} \emph{machchhar ke kāṭne par voh hāth khujā rahā thā.} — He was scratching his hand where the mosquito had bitten.
  \item \emph{Recall:} say the Urdu for perfume, then the Urdu for the sense of touch, then say \emph{khujānā} again
\end{itemize}

\begin{tcolorbox}[breakable,colback=teal!4,colframe=teal!35!black,title={Before you move on}]
What does \ur{کھجانا} mean? (\textquotedblleft{}To scratch an itch\textquotedblright{}.) Say it once more.
\end{tcolorbox}
\section[\texorpdfstring{\ur{خارش} (khārish)}{خارش (khārish)}]{\ur{خارش} (khārish) — a rash; an itch (on the skin)}

\label{lesson:UR-C468-kharish}



\begin{quote}
Before the new one: say the Urdu for the sense of touch, then the Urdu for to scratch an itch.
\end{quote}

\subsection*{You'll want to know: \ur{خارش}}
\textbf{\ur{خارش}} — \emph{khārish} — \textquotedblleft{}a rash; an itch (on the skin)\textquotedblright{}.

\begin{quote}
\ur{اس} \ur{کی} \ur{جلد} \ur{پر} \ur{خارش} \ur{ہے۔}
\end{quote}

\begin{quote}
\emph{us kī jild par khārish hai.}
\end{quote}

\begin{quote}
She has a rash on her skin.
\end{quote}

Feminine.

\begin{cousinweb}
\textbf{\ur{خارش}} (\emph{khārish}) is usually traced to Persian \emph{khārish}, \textquotedblleft{}itching\textquotedblright{}.
\end{cousinweb}

\begin{grammarlens}[title={What you've built}]
One more word for describing what you lived and felt.
\end{grammarlens}

\subsection*{Your turn}
\begin{itemize}
\raggedright
  \item \emph{Say it:} \emph{khārish}
  \item \emph{Say it:} \emph{khārish}, once more
  \item \emph{Say it:} \emph{us kī jild par khārish hai.} — She has a rash on her skin.
  \item \emph{Recall:} say the Urdu for the sense of touch, then the Urdu for to scratch an itch, then say \emph{khārish} again
\end{itemize}

\begin{tcolorbox}[breakable,colback=teal!4,colframe=teal!35!black,title={Before you move on}]
What does \ur{خارش} mean? (\textquotedblleft{}A rash; an itch (on the skin)\textquotedblright{}.) Say it once more.
\end{tcolorbox}
